\chapter{知识的整合}
\label{chap:ket-integration}

花卉识别系统已有两个分别训练的分类器。模型一主要见过清晰的近景照片，模型二见过更多低清晰度照片；两者识别同一组物种，却会在不同照片上出错。系统还保存了识别颜色、形态和病斑的适配模块。选取一个来源会放弃其他来源的判断，但把它们全部放进系统也还没有规定如何共同使用。我们需要比较三种做法：保留来源并组合计算，读取来源参数形成新模型，或者用来源的回答训练另一个模型。怎样选择，取决于希望保留什么能力、能访问什么信息，以及部署允许付出多少代价。

将分别获得的模型、模块或其他载体共同用于所需行为，称为\term{知识整合}（knowledge integration）。这里沿用概述章对知识的操作性含义：参数和表示提供支持预测与学习的信息，不必能分解成逐条语义事实。本章先固定花卉物种分类的输入、类别含义和独立评价样本，比较来源进入结果的不同方式；转向属性判断、指令或动作组合时，再说明新增要求。所有数值算例都是教学构造，不是实验测量。

图~\ref{fig:ket-integration-routes}的三个面板分别标出来源参与的计算。\term{模型组合}（model composition）保留来源单元，让它们在预测时按约定共同运行。\term{参数融合}（model merging）直接对来源权重实施运算，随后运行构造出的模型。\term{知识蒸馏}（knowledge distillation）则让来源的输出或表示充当监督，训练承接模型。观察图中的部署单元，可以看出组合仍需调用来源，后两种路线不再逐次调用原来源。这些机制可以独立采用，也可以串联，例如先组合两个来源，再用组合行为训练承接模型；它们不是三个必经阶段。表~\ref{tab:ket-integration-route-choice}按访问时机、额外训练和部署状态比较这些选择，准备与运行阶段的成本不能混为一项。

\begin{table*}[htbp]
  \centering
  \small
  \begin{tabularx}{\linewidth}{@{}p{22mm}XXXX@{}}
    \toprule
    路线 & 来源访问 & 额外材料与训练 & 部署形态 & 首要失败边界 \\
    \midrule
    模型组合 & 部署时能够执行来源；模块连接还需中间接口 & 固定规则可直接使用；学习聚合或路由需要相应材料 & 保留所需来源、聚合器或路由器 & 接口错位、共同错误、错误路由与重复计算 \\
    参数融合 & 准备时能够读取并对应来源参数及状态 & 可只用参数，也可用输入统计、激活或标签决定运算 & 通常保存一份参数，也可能保留任务分支 & 坐标不对应、平均经过高损失区域与任务冲突 \\
    知识蒸馏 & 训练时能够取得教师响应；表示约束还需内部激活 & 需要查询输入、学生训练和独立模型选择 & 通常只保留学生的预测路径 & 输入覆盖不足、教师偏差与学生容量不足 \\
    \bottomrule
  \end{tabularx}
  \caption{三种整合方式的访问、训练和部署条件。各行比较的是来源如何参与结果，条件可以交叉，例如带路由的组合系统也能充当蒸馏教师。部署只保留一个模型，不表示准备阶段无需数据、教师存储或额外计算。}
  \label{tab:ket-integration-route-choice}
\end{table*}

\begin{figure*}[htbp]
  \centering
  % Three independent integration routes; schematic, no empirical data.
% Final canvas: 164 mm x 131 mm. Double outlines mark deployed units.
% Solid: forward computation; dashed: training; dotted: parameter reading.
\begin{tikzpicture}[
  x=1mm,y=1mm,font=\figurefont\footnotesize,text=BookInk,
  box/.style={rectangle,draw=FigureModel,fill=FigureModel!7,
    line width=0.8pt,minimum width=25mm,minimum height=9mm,
    inner sep=1mm,align=center},
  deployed/.style={box,double,double distance=0.7pt},
  input/.style={box,draw=FigureInput,fill=FigureInput!7,
    minimum width=15mm},
  result/.style={box,draw=FigureOutput,fill=FigureOutput!7,
    minimum width=15mm},
  loss/.style={box,draw=FigureLoss,fill=FigureLoss!7},
  flow/.style={-{Latex[length=2mm,width=1.3mm]},draw=BookInk,
    line width=0.8pt,rounded corners=1mm},
  train/.style={flow,dashed,draw=FigureLoss},
  weights/.style={flow,dotted,draw=BookBlue},
  construct/.style={flow,draw=BookBlue},
  note/.style={font=\figurefont\footnotesize,text=BookMuted,inner sep=1mm}
]
  \path[use as bounding box] (0,-13) rectangle (164,118);
  \node[anchor=west,font=\figurefont\small] at (1,113) {(a) 保留来源计算};
  \node[input] (xa) at (12,94) {$\bm{x}$};
  \node[deployed] (a1) at (47,101) {来源$f_1$};
  \node[deployed] (a2) at (47,87) {来源$f_2$};
  \node[deployed] (agg) at (99,94) {聚合器};
  \node[result] (ya) at (150,94) {预测};
  \draw[flow] (xa.east) -- (25,94) |- (a1.west);
  \draw[flow] (25,94) |- (a2.west);
  \draw[flow] (a1.east) -- (76,101) |- (agg.west);
  \draw[flow] (a2.east) -- (76,87) |- (agg.west);
  \draw[flow] (agg.east) -- (ya.west);
  \draw[BookRule,line width=0.6pt] (0,78)--(164,78);

  \node[anchor=west,font=\figurefont\small] at (1,73) {(b) 读取来源参数};
  \node[box] (b1) at (34,60) {来源$f_1$\\$\bm{\theta}_1$};
  \node[box] (b2) at (34,45) {来源$f_2$\\$\bm{\theta}_2$};
  \node[box] (rule) at (76,53) {融合规则};
  \node[deployed] (merged) at (122,53) {融合模型};
  \node[result] (yb) at (151,53) {预测};
  \node[input] (xb) at (92,37) {$\bm{x}$};
  \draw[weights] (b1.east) -- (53,60) |- (rule.west);
  \draw[weights] (b2.east) -- (53,45) |- (rule.west);
  \draw[construct] (rule.east) -- node[note,above] {构造新参数} (merged.west);
  \draw[flow] (xb.east) -| (merged.south);
  \draw[flow] (merged.east) -- (yb.west);
  \draw[BookRule,line width=0.6pt] (0,29)--(164,29);

  \node[anchor=west,font=\figurefont\small] at (1,24) {(c) 用来源监督训练};
  \node[input] (xc) at (12,8) {$\bm{x}$};
  \node[box] (c1) at (45,15) {来源$f_1$};
  \node[box] (c2) at (45,1) {来源$f_2$};
  \node[loss] (kd) at (85,8) {蒸馏约束};
  \node[deployed] (student) at (124,8) {学生};
  \node[result] (yc) at (152,8) {预测};
  \draw[flow] (xc.east) -- (24,8) |- (c1.west);
  \draw[flow] (24,8) |- (c2.west);
  \draw[train] (c1.east) -- (65,15) |- (kd.west);
  \draw[train] (c2.east) -- (65,1) |- (kd.west);
  \draw[train] (kd.east) -- node[note,above] {更新} (student.west);
  \draw[flow] (student.east) -- (yc.west);
  \draw[flow] (student.north) -- (124,21) -| (kd.north);
  \node[note] at (105,23) {学生输出};
  \draw[flow] (xc.south) -- (12,-9) -| (student.south);
  \node[note,fill=white] at (86,-9) {同一输入};
\end{tikzpicture}

  \caption{同一对来源模型参与整合的三种方式。实线表示前向计算，虚线表示监督约束及学生更新，点线表示参数读取，蓝色实线表示融合规则生成新参数。双重轮廓标出部署时仍需保存并运行的单元：组合需要来源与聚合器，融合运行新模型，蒸馏运行学生。三个面板互相独立，不表示操作先后。}
  \label{fig:ket-integration-routes}
\end{figure*}

来源是否可用、相对不复用是否有收益，沿用第~\ref{chap:ket-reuse-selection}~章的判断；改变系统后怎样保留能力和管理版本，沿用第~\ref{chap:ket-update-maintenance}~章的要求。本章在这些条件下解释现成来源如何共同发挥作用。多个任务共同塑造共享计算则由第~\ref{sec:ket-sharing-computation}~节展开。整合也可能训练组合器或学生，因而不能以“有没有训练”区分两者；要看当前利用的是已形成的来源，还是多个任务正在共同形成知识。训练材料、标签与记录的产生和校验仍归知识获取主题，即第~\ref{part:efficient-knowledge-use}。

\section{基于模型组合的知识整合}
\label{sec:ket-integration-composition}

两个来源保留原参数以后，仍需约定信息在哪里汇合。若两者都给出物种分布，可以汇总对同一个问题的预测；若前一模块给出花瓣属性，后一模块据此作判断，就要连接中间计算。两种做法分别要求答案含义可比和上下游接口相容。完整网络与小型适配器都可能承担这些角色，模块大小并不决定属于哪种组合。

组合是否改变，还要相对于预先声明的服务范围判断。在该任务与输入范围内，所用来源、系数和实际连接保持不变，就是固定组合；依据任务或当前输入选择来源、权重或结构，就是条件组合。通过训练学得后再冻结的全局权重仍是固定组合，人工规定的输入路由仍是条件组合。这里区分的是运行规则对条件的依赖，不是规则有没有经过训练。

\subsection{预测结果的集成}
\label{sec:ket-integration-prediction}

在花卉分类任务中，先把所有来源的输出对应到同一个类别集合$Y=\{1,\ldots,c\}$，其中$c$为类别数。第$k$个来源在照片$\bm{x}$上给类别$y$的概率记为$p_k(y\mid\bm{x})$。若两个向量分别按“月季、百合、菊花”和“百合、菊花、月季”排列，同一坐标就不是同一物种，必须先重排。物种与病害输出、多标签概率与互斥类别概率也不能仅凭向量等长就相加。

类别对应正确以后，还要检查数值是否具有可比含义。一个来源报告置信度$0.8$时，相关样本的实际正确率可能明显低于$0.8$；检查这类概率与观测频率的关系，就是校准检查。它需要与拟服务范围相符的有标签材料，不能只看数值是否归一化。\BookChapterRef{chap:ensemble-learning}{卷4《模型：经典模型》第6章“集成学习”}已讨论成员的产生与训练，这里把来源视为已选定，进一步比较怎样使用答案和承担运行成本。

\subsubsection{基于固定规则的预测集成}
\label{sec:ket-integration-fixed-prediction}

对同一张照片，一个来源偏向月季，另一个偏向百合，怎样作出最终判断？多数投票只计每个来源选择的类别，再选票数最多的一类；两票各不相同时需要预先规定平局处理。若还希望利用来源对其他类别的倾向，可以保留概率向量，按固定权重作算术平均：
\begin{equation}
  \bar p(y\mid\bm{x})=\sum_{k=1}^{K}\alpha_k p_k(y\mid\bm{x}),
  \qquad \alpha_k\geq 0,\quad \sum_{k=1}^{K}\alpha_k=1.
  \label{eq:ket-integration-fixed-average}
\end{equation}
其中$K$为来源数，$\alpha_k$为第$k$个来源的权重，在声明范围内不随输入变化。非负权重避免出现负概率，权重和为一使各类概率之和仍为一；这只保证分布合法，不保证校准或分类正确。若需从数据确定$\bm{\alpha}=(\alpha_1,\ldots,\alpha_K)^\top$，可以在与来源拟合隔离的验证材料上最小化约定损失，再冻结权重。最终测试还应使用未参与权重选择的材料。

\begin{example}[互补错误与共同错误]
  \label{ex:ket-integration-weighted-prediction}
  将花卉任务简化为月季与非月季二分类。正类$y=1$表示月季，正类概率达到$0.5$时判为月季，平局也归正类。四张教学照片的真实标签与来源概率为
  \[
    (y_1,y_2,y_3,y_4)=(1,0,1,0),\qquad
    \begin{aligned}
      (p_1(\bm{x}_i))_{i=1}^{4}&=(0.9,0.4,0.6,0.6),\\
      (p_2(\bm{x}_i))_{i=1}^{4}&=(0.4,0.1,0.9,0.3).
    \end{aligned}
  \]
  这里$p_k(\bm{x}_i)$简写$p_k(y=1\mid\bm{x}_i)$。模型一错第四张，模型二错第一张，两者正确率都是$3/4$。给两个来源权重$0.6$与$0.4$，组合概率为
  \[
    (\bar p(\bm{x}_i))_{i=1}^{4}
    =(0.70,0.28,0.72,0.48).
  \]
  第四张由$0.6\times0.6+0.4\times0.3=0.48$落到阈值以下，四张照片因而全部分类正确。现在保持标签、模型一和权重不变，只将模型二概率换成$(0.8,0.2,0.7,0.7)$。模型二正确率仍是$3/4$，但它与模型一都错第四张；组合概率变为$(0.86,0.32,0.64,0.64)$，第四张仍错。成员数量和各自正确率相同，错误发生在哪些输入上却不同，组合收益也随之改变。
\end{example}

算例指出错误依赖的重要性，但阈值分类不便直接拆解。为了看清平均改变了哪一部分误差，我们暂时改用实值回归和平方损失。设$y=f_\star(\bm{x})+\varepsilon$，其中$f_\star(\bm{x})=\mathbb{E}[y\mid\bm{x}]$，噪声满足$\mathbb{E}[\varepsilon\mid\bm{x}]=0$。固定输入$\bm{x}$，只让训练样本和训练随机性变化，假设所需二阶矩有限，并假设测试噪声与训练过程在该输入下条件独立。

将第$k$个预测器相对回归函数的误差记为$e_k=\hat f_k(\bm{x})-f_\star(\bm{x})$，误差均值为$b_k=\mathbb{E}[e_k]$，协方差为$\Sigma_{k,j}=\operatorname{Cov}(e_k,e_j)$。这些期望描述重复训练的波动，不是对当前几个固定预测值取平均。固定权重的和为一时，组合误差就是$\bar e=\sum_k\alpha_ke_k$。把均值和中心化误差分开，可以得到它的期望平方：
\begin{equation}
  \begin{aligned}
    \mathbb{E}[\bar e^2]
    &=\biggl(\sum_k\alpha_k b_k\biggr)^2+
      \mathbb{E}\biggl[\biggl(\sum_k\alpha_k(e_k-b_k)\biggr)^2\biggr]\\
    &=\biggl(\sum_k\alpha_k b_k\biggr)^2+
      \sum_k\alpha_k^2\Sigma_{k,k}
      +2\sum_{k<j}\alpha_k\alpha_j\Sigma_{k,j}.
  \end{aligned}
  \label{eq:ket-integration-bias-covariance}
\end{equation}
第一行的交叉项期望为零，因为中心化误差均值为零；第二行展开方差，每对不同来源贡献一个协方差项。公式因此既保留组合偏差，也保留来源错误一起变化的影响。若计分对象是带噪声标签$y$，还需加上$\operatorname{Var}(\varepsilon\mid\bm{x})$，测试噪声的零均值与条件独立使噪声交叉项消失；若评价整个目标输入范围，则再对其输入分布取期望。

用一个对称设定观察来源数量的作用：各来源无偏，方差都为$\sigma^2>0$，任意两个来源的相关系数都为$\rho$，并取等权平均。将这些值代入协方差展开，得到
\begin{equation}
  \operatorname{Var}(\bar e)
  =\frac{K\sigma^2+K(K-1)\rho\sigma^2}{K^2}
  =\sigma^2\left(\rho+\frac{1-\rho}{K}\right).
  \label{eq:ket-integration-equal-correlation}
\end{equation}
完全相关时$\rho=1$，平均没有降低方差；不相关时$\rho=0$，方差降到$\sigma^2/K$。等相关协方差必须有效，$K>1$时应有$-1/(K-1)\leq\rho\leq1$。若来源还有相同的非零偏差，平均也不会去掉它。回到分类任务，概率是否越过决策阈值还会影响对错，不能把这条平方损失关系当作准确率必然提高的证明。

分类中也要固定汇总的位置。来源未归一化分数为$\bm{a}_k$，经$\operatorname{softmax}$得到概率，一般有$\operatorname{softmax}(\sum_k\alpha_k\bm{a}_k)\neq\sum_k\alpha_k\operatorname{softmax}(\bm{a}_k)$。在当前归一权重下，先平均分数再归一化，相当于把严格正的类别概率作加权几何平均后归一化；直接平均概率则作算术平均。前者更容易受分数尺度影响，投票又丢弃了概率幅度。未经校准的高置信度不能直接作为提高来源权重的依据。

这些固定规则通常需要执行所有非零权重的来源。独立模型的权重存储和运算量大致相加；并行执行可以减少等待时间，却需要设备并受到通信与最慢来源的限制，串行执行则累积延迟。共享底座的适配器可能共用部分计算，但应按实际路径核算。是否节约资源和是否降低任务错误，是两项需要分别验证的结论。

\subsubsection{基于条件规则的预测集成}
\label{sec:ket-integration-conditional-prediction}

若模型一在清晰照片上可靠、模型二在模糊照片上可靠，全局权重就无法按这种差别选择。条件规则让选择器读取预测时可得的图像质量、任务说明或其他证据，再分配权重。用$u$索引目标任务$\mathcal T_u$，与来源索引$k$分开；$\bm{q}(\bm{x},u)$表示选择器可观察的信息。同一任务内，只有输出含义相容且可执行的来源进入候选，组合分布写为
\begin{equation}
  \bar p(y\mid\bm{x},u)
  =\sum_k\alpha_k(\bm{q}(\bm{x},u))p_k(y\mid\bm{x},u).
  \label{eq:ket-integration-conditional-average}
\end{equation}
若要将结果解释为混合分布，权重仍应非负且和为一。只保留一个来源是硬选择，多个来源共同贡献是软组合。学习选择器时，可在与来源拟合隔离的有标签材料上估计不同质量区域的损失，也可直接训练门控函数以减小组合损失。门控函数就是从可见条件到权重的映射。其训练材料应覆盖拟服务条件，不能用来源在自身训练样本上的过拟合成绩证明条件可靠性。

仍以月季二分类说明路由的得失，阈值沿用$0.5$。一张清晰月季和一张模糊非月季上，模型一概率为$0.9,0.95$，模型二为$0.2,0.2$。等权平均得到$0.55,0.575$，错第二张。若清晰度指标正确区分两张照片，清晰照片调用模型一、模糊照片调用模型二，结果为$0.9,0.2$，均分类正确。若两张照片的质量分组恰好颠倒，则得到$0.2,0.95$，两张都错。选择依据是可见质量，不是预测时未知的真实答案。

条件发生在任务层还是输入层，会改变调用成本。按任务说明或少量示例确定一次$\bm{\alpha}_u$，然后在该任务内冻结，是任务间的条件组合；逐张照片选择，则每次都要执行选择器。从原图低成本估计清晰度，可能只需运行选中的来源；若必须先取得全部模型的置信度，来源计算已经发生，最后只保留一个答案不会撤销这些费用。

对超出已验证条件或缺少选择证据的输入，可以预设弃权、人工复核或备用流程。评价弃权时，要同时报告被接纳输入占全部输入的比例，以及接纳部分的错误率；只接纳极少容易样本不能表示系统整体准确。地区、季节和设备变化还可能改变质量指标与来源误差的关系，需要重新检验路由。这里组合的是已有答案，对答案加权并不等于组合来源内部的推理步骤。

\subsection{计算模块的组合}
\label{sec:ket-integration-modules}

物种概率可以回答“这是什么花”，却未必能回答“图中哪朵花同时呈红色和钟形”。后一目标需要保留属性及其所属对象，再继续作判断，因而从同义答案的汇总转向计算模块的连接。对单花图像，可以传递属性向量；对多花图像，还要把颜色、形态与位置绑定到同一对象，避免取一朵花的颜色和另一朵花的形态。图~\ref{fig:ket-integration-interfaces}对照了这两种接口：左侧汇总分布，右侧让中间表示继续参与计算。

\begin{figure*}[htbp]
  \centering
  % Interface comparison, two fixed tasks; schematic, no empirical data.
% Canvas: 164 mm x 69 mm. Dashed arrows carry control, not predictions.
\begin{tikzpicture}[
  x=1mm,y=1mm,font=\figurefont\footnotesize,text=BookInk,
  box/.style={rectangle,draw=FigureModel,fill=FigureModel!7,
    line width=0.8pt,minimum width=22mm,minimum height=9mm,
    align=center,inner sep=1mm},
  input/.style={box,draw=FigureInput,fill=FigureInput!7,
    minimum width=15mm},
  result/.style={box,draw=FigureOutput,fill=FigureOutput!7,
    minimum width=20mm},
  control/.style={box,draw=BookAmber,fill=BookAmber!7},
  flow/.style={-{Latex[length=2mm,width=1.3mm]},draw=BookInk,
    line width=0.8pt,rounded corners=1mm},
  choose/.style={flow,dashed,draw=BookAmber},
  edge label/.style={fill=white,inner sep=0.6mm,
    font=\figurefont\footnotesize}
]
  \path[use as bounding box] (0,0) rectangle (164,69);
  \node[anchor=west,font=\figurefont\small] at (0,65) {(a) 同义预测的集成};
  \node[anchor=west,text=BookMuted] at (0,58) {任务：识别物种};
  \node[input] (image-a) at (9,33) {图像};
  \node[box] (model-a) at (39,45) {来源$f_1$};
  \node[box] (model-b) at (39,21) {来源$f_2$};
  \node[result] (aggregate) at (68,33) {聚合};
  \node[result] (species) at (68,6) {物种};
  \draw[flow] (image-a.east) -- (22,33) |- (model-a.west);
  \draw[flow] (22,33) |- (model-b.west);
  \draw[flow] (model-a.east) -- (68,45) -- (aggregate.north);
  \draw[flow] (model-b.east) -- (68,21) -- (aggregate.south);
  \node[edge label] at (65,51) {$p_1(y\mid\bm{x})$};
  \node[edge label] at (64,15) {$p_2(y\mid\bm{x})$};
  \draw[flow] (aggregate.east) -- (80,33) |- (species.east);
  \draw[BookRule,line width=0.6pt] (85,1)--(85,68);

  \node[anchor=west,font=\figurefont\small] at (90,65) {(b) 中间计算的组合};
  \node[anchor=west,text=BookMuted] at (90,58) {任务：属性条件判断};
  \node[input] (image-b) at (101,43) {图像};
  \node[box] (attributes) at (144,43) {属性模块$g$};
  \node[control] (selector) at (104,21) {条件选择器};
  \node[box] (reasoner) at (144,21) {判断模块$h$};
  \node[result] (decision) at (144,5) {判断};
  \draw[flow] (image-b.east) -- (attributes.west);
  \draw[flow] (attributes.south) -- (reasoner.north);
  \node[edge label,anchor=west] at (147,32) {$\bm{z}$};
  \node[edge label,anchor=east] at (142,32) {中间表示};
  \draw[flow] (image-b.south) -- (101,32) -| (selector.north);
  \draw[choose] (selector.east) -- node[edge label,above] {选择信号} (reasoner.west);
  \draw[choose] (selector.west) -- (89,21) -- (89,52)
    -| (attributes.north);
  \draw[flow] (reasoner.south) -- (decision.north);
\end{tikzpicture}

  \caption{答案汇总与中间计算连接的差别。左侧固定物种分类，两条来源接口携带类别顺序相同的分布；右侧固定属性条件判断，属性模块向判断模块传递中间表示，虚线仅传递模块选择信号。单花示例中，若上游按“红色、钟形”输出，下游却按反序读取，两维接口虽然能连接，实际判断已改变。多花场景还需额外保存属性所属对象。}
  \label{fig:ket-integration-interfaces}
\end{figure*}

\subsubsection{基于固定连接的模块组合}
\label{sec:ket-integration-fixed-modules}

先把输入限定为只有一朵花的照片，目标是判断它是否同时满足两种属性。属性模块$g$输出$\bm{z}=g(\bm{x})\in[0,1]^2$，依次表示红色与钟形的置信度；固定判断器采用教学规则$h(\bm{z})=\mathbf{1}\{z_1>0.8,\ z_2>0.6\}$，指示函数在两条件都满足时返回真。输出$(0.9,0.7)^\top$经过$h$会返回真。若下游反序读取为$(0.7,0.9)^\top$，红色阈值不再满足；若顺序正确却把钟形错估成$0.4$，也会返回假。结果相同，前一种是接口错误，后一种是属性估计错误。

这种固定连接在声明范围内始终计算$h(g(\bm{x}))$。将颜色和形态分交两个来源，也可拼接后计算$h([g_{\mathrm{color}}(\bm{x});g_{\mathrm{shape}}(\bm{x})])$，其中分号表示向量拼接。下游仍须实现“两条件同时成立”的关系，拼接本身不会完成判断。真实接口还要约定单位、归一化尺度、位置、缺失值和状态；视频模块若需要连续帧及历史缓存，就不能因为单帧张量形状相同而替换输入。

当上游输出与下游熟悉的表示不一致时，可以在两者之间学习映射$r_{\bm{\phi}}$，组成$h(r_{\bm{\phi}}(g(\bm{x})))$。若每个训练输入都有对应的参考表示$\bm{z}^{\mathrm{ref}}_i$，平方残差$\sum_i\lVert r_{\bm{\phi}}(g(\bm{x}_i))-\bm{z}^{\mathrm{ref}}_i\rVert_2^2$就直接约束表示对应。若只有最终标签，则固定来源后用整链任务损失学习$\bm{\phi}$，但需要可微路径或另一个可行求解方法。最终标签只约束最终判断，不能保证映射的每一维恢复了预定属性含义；是否还更新上下游参数也要另行声明。

即使语义接口正确，下游也可能不习惯预测中间值。它若只用真实属性的零一标签训练，接入上游后面对的却是软概率、缺测与相关错误。分别用真实中间值和来源预测测试下游，可以区分这项输入变化与下游判断本身的错误；也可将来源预测用于下游训练，再用独立样本检查整链。冻结来源参数只维持原输入、状态和调用方式下的函数，不保护更改后的输入分布、路由和最终决策。

\subsubsection{基于条件选择的模块组合}
\label{sec:ket-integration-conditional-modules}

判断红色钟形花和判断病斑叶片，可能需要不同提取器与连接。搜索之前要限定允许使用的模块和合法接口，避免把任意连线都当作候选。\term{组合语法}（composition grammar）规定模块可出现的位置及允许的输入输出连接。\term{模块化元学习}（modular meta-learning）先在相关任务中学习可重组模块，再用新任务的少量样本选择结构。本节固定已形成的模块库，考察后一步；任务分布的准备与留出条件分别由第~\ref{sec:ket-transfer-preparation}~节和第~\ref{sec:ket-transfer-task-splits}~节展开\scite{ket-alet2018-modular}

用四个教学模块说明选择。$g_1,g_2$提取属性，$h_1,h_2$接受相容格式并作判断，组合语法只允许一个提取器接一个判断器。因此候选集合是$\mathcal{G}=\{h_b\circ g_a:a,b\in\{1,2\}\}$，$\circ$表示函数复合。新任务提供非空支持集$S_u$，即供任务内选择使用的带标签样本。按单样本损失$\ell$比较四个结构，选择平均损失最小的候选：
\begin{equation}
  G_u^\star\in\argmin_{G\in\mathcal{G}}
  \frac{1}{|S_u|}\sum_{(\bm{x},y)\in S_u}\ell(G(\bm{x}),y).
  \label{eq:ket-integration-structure-search}
\end{equation}
模块参数在搜索中不变，改变的是计算结构；有限且非空的候选集使最小值可取到，并列时可预先指定处理规则。若按$h_1\circ g_1,h_2\circ g_1,h_1\circ g_2,h_2\circ g_2$排列的四个候选，在四个支持样本上分别错$2,1,3,0$个，零一损失的均值为$0.5,0.25,0.75,0$，就选第四个。还需用独立查询集$Q_u$检验未见输入，支持集上的零错误不保证任务风险为零。

结构多时，穷举需要执行过多候选。模块化元学习可从一个合法结构出发，替换模块或改变合法连接，并用模拟退火接受候选：损失不增加时接受，增加时也以一定概率接受，避免每一步都被局部改善限制。相同温度下，损失增量越大，接受概率越低；对同一正增量，降温会降低接受概率。四模块例子可以依次替换判断器与提取器，但有限预算的搜索未必找到全局最优。任务内固定所选结构时只执行相应模块；逐输入改变结构时，还要承担在线路由与切换。

也可以保留连接，在每层改变来源表示的贡献。AdapterFusion冻结底座和已训练适配器，学习融合层，让当前表示决定如何使用各适配器输出。它无需先选出唯一模块，而是将候选表示投影后加权，供后续网络继续计算。设底座表示$\bm{h}\in\mathbb{R}^{d_h}$、适配器输出$\bm{z}_k\in\mathbb{R}^{d}$，查询与键投到$d_a$维，值投到$d_v$维，其核心注意力可概括为
\begin{equation}
  \begin{aligned}
    \bm{q}&=\bm{W}_{Q}\bm{h},&
    \bm{k}_k&=\bm{W}_{K}\bm{z}_k,&
    \bm{v}_k&=\bm{W}_{V}\bm{z}_k,\\
    a_k&=\frac{\exp(\bm{q}^{\top}\bm{k}_k)}
                  {\sum_j\exp(\bm{q}^{\top}\bm{k}_j)},&
    \bm{z}_{\mathrm{fuse}}&=\sum_k a_k\bm{v}_k.
  \end{aligned}
  \label{eq:ket-integration-adapter-attention}
\end{equation}
这里$\bm{W}_Q\in\mathbb{R}^{d_a\times d_h}$、$\bm{W}_K\in\mathbb{R}^{d_a\times d}$、$\bm{W}_V\in\mathbb{R}^{d_v\times d}$，$\bm{q},\bm{k}_k,\bm{v}_k$分别是查询、键和值。查询与键的内积给来源计分，指数归一化得到非负且和为一的$a_k$，再加权值向量。融合输出为$d_v$维，须与后续接口相容；此处省略层、位置与残差细节。目标标签训练这些投影及相应任务头，来源模块不更新，因此发生的是激活组合，不是离线参数平均。软注意力仍需取得所有候选表示，通常不能省掉其计算。原论文主要实验包含目标任务已训练的适配器，不能据此声称完全未见任务也满足同样条件\scite{ket-pfeiffer2021adapterfusion}

与第~\ref{sec:ket-sharing-conditional}~节的多门控混合专家相比，AdapterFusion固定来源，后者还由多个任务共同训练专家。形式相似的权重不能抹去这种训练差异。结构选择或表示融合能否带来组合泛化，也要另测：组合泛化是用见过的成分完成未见组合，而不只是回答来源任务。属性测试可保留红色与钟形的单独出现，把两者同时成立留出；指令测试可分别训练筛选和排序，留出先筛选后排序。应比较单一来源、固定组合与重组，并检查中间结果是否真正被使用。Conklin等用相似性驱动的任务采样形成训练与推广任务对，为这种能力提供训练安排；它不保证任意合法模块组合都能推广\scite{ket-conklin2021compositional}

从静态图像转向技能时，接口还包含执行后的环境状态。花卉采样机器人先接近植株、再对准相机，前一技能结束的状态必须属于后一技能可处理的范围。另一种整合不是顺序串接，而是综合已有策略对动作的价值判断。\term{后继特征}（successor features）在固定策略下记录未来特征的折扣累计期望；策略决定各状态如何选择动作，折扣使较远的特征贡献较小。若动力学、特征和状态动作接口共享，且奖励是这些特征的线性函数，后继特征与新奖励权重的内积就能评价旧策略在新任务的价值。广义策略改进再比较各旧策略给出的动作价值来选择动作。价值整合不保证任意技能终止状态可以衔接，其完整推导可参考\BookVolumeRef{vol:reinforcement-learning}{卷7《学习范式：强化学习》}\scite{ket-barreto2016successor}

模块库增长以后，结构候选还依赖版本与检索。冗余模块增加搜索量，删除模块会影响引用它的组合，漏检则使本来可用的模块没有参加选择。应保存组合所用版本、接口、依赖和回退清单，分别计入检索、搜索、加载与执行成本。第~\ref{sec:ket-update-versions}~节的版本要求和第~\ref{sec:ket-overview-information}~节的访问要求仍适用；本节增加的检查是这些已知单元在指定路径上能否共同完成任务。

\section{基于参数融合的知识整合}
\label{sec:ket-integration-merging}

花卉终端若只能保存并运行一份模型，多来源组合就可能超出部署预算。参数融合改为在准备阶段读取来源权重，直接生成新的参数，再用它完成预测。这要求先确认哪些坐标可以一起运算：数组形状相同，不表示同一位置承担相同计算。对应确定后，才能选择对完整权重、共享层或相对底座增量的融合规则，再用允许的信息确定系数。减少部署模型数并不保证保留了全部来源能力。

\subsection{参数对应关系的建立}
\label{sec:ket-integration-alignment}

要看清坐标为何影响平均，可以先交换一个网络的隐藏单元，而不改变它的函数。设输入$\bm{x}\in\mathbb{R}^d$，隐藏层有$m$个单元，输出为$c$维，激活$\sigma$逐元素作用：
\[
  f(\bm{x})=\bm{W}_2\sigma(\bm{W}_1\bm{x}+\bm{b}_1)+\bm{b}_2,
  \quad \bm{W}_1\in\mathbb{R}^{m\times d},\
  \bm{W}_2\in\mathbb{R}^{c\times m}.
\]
偏置为$\bm{b}_1\in\mathbb{R}^m$、$\bm{b}_2\in\mathbb{R}^c$。置换矩阵$\bm{P}\in\mathbb{R}^{m\times m}$仅重排隐藏单元，满足$\bm{P}^{\top}\bm{P}=\bm{I}$；逐元素激活也与这种重排可交换，即$\sigma(\bm{P}\bm{v})=\bm{P}\sigma(\bm{v})$。为了保持每个隐藏输出仍连到原来的输出权重，应同时变换两层：
\begin{equation}
  \begin{aligned}
    \widetilde{\bm{W}}_1&=\bm{P}\bm{W}_1,&
    \widetilde{\bm{b}}_1&=\bm{P}\bm{b}_1,\\
    \widetilde{\bm{W}}_2&=\bm{W}_2\bm{P}^{\top},&
    \widetilde{\bm{b}}_2&=\bm{b}_2
  \end{aligned}
  \label{eq:ket-integration-permutation}
\end{equation}
代回网络，输出为$\bm{W}_2\bm{P}^{\top}\bm{P}\sigma(\bm{W}_1\bm{x}+\bm{b}_1)+\bm{b}_2=f(\bm{x})$。隐藏层权重换行，后一层同步换列，因而中间坐标变了，函数却没有变。只改一层则通常不能抵消置换。这个等价关系正好说明：函数本身不能唯一确定可逐坐标平均的参数数组。

\begin{example}[完全相同的函数也不能直接平均]
  \label{ex:ket-integration-permutation}
  取一个标量输入、标量输出的教学网络，偏置为零，$\operatorname{ReLU}(v)=\max(v,0)$。两层权重和隐藏单元交换为
  \[
    \bm{W}_1=\begin{bmatrix}1\\-1\end{bmatrix},\quad
    \bm{W}_2=\begin{bmatrix}1&1\end{bmatrix},\quad
    \bm{P}=\begin{bmatrix}0&1\\1&0\end{bmatrix}.
  \]
  两个来源都计算$\operatorname{ReLU}(x)+\operatorname{ReLU}(-x)=|x|$。直接平均第一层的$(1,-1)^\top$与$(-1,1)^\top$，却得到零向量，后一层不变时输出恒为零。于是$x=1$上两个来源均输出1，融合结果输出0。先对第二个来源实施逆置换，两份数组恢复一致，平均才仍计算$|x|$。失败不是来源函数不同，而是平均前没有恢复对应。
\end{example}

Git Re-Basin搜索这种隐藏单元对应。激活匹配在共同输入上给单元对计算相似度，再求一一指派；权重匹配只读取参数，逐层更新置换以增大对应权重的匹配程度。一一指派要求每个单元只配给一个单元，邻层与残差连接还须配套重排。前者需要代表性输入，后者无需读取训练样本，但匹配目标都不是全部插值模型的任务风险。所谓低损失线性连通，是指对应后的参数线段沿途没有明显损失障碍，不能仅从找到一个置换推出。论文还构造了不能经任意置换获得这种连通的非随机梯度下降解；它限制的是对任意解的推广，并不反驳针对随机梯度下降所得解的限定猜想\scite{ket-ainsworth2023rebasin}

若来源分别识别物种和病害，一一匹配也可能把用途不同的特征硬配成对。ZipIt!用无标签输入上的激活相关性，不仅匹配跨来源特征，也合并同一来源内部的冗余特征；还可只融合前面若干层，保留后续任务分支。这改变了匹配限制，不只是重新排列模型二的单元。保留分支后，部署仍可能需要任务身份和相应预测头，计算与存储应按这张实际结构核算\scite{ket-stoica2024zipit}

低秩更新的因子还会引入内部坐标。\term{低秩适应}（low-rank adaptation, LoRA）以两个矩阵的乘积表示相对底座的权重变化，用较小的内维数限制增量的秩。某层写为$\Delta\bm{W}_k=\bm{B}_k\bm{A}_k$，其中$\bm{B}_k\in\mathbb{R}^{d_{\mathrm{out}}\times r}$、$\bm{A}_k\in\mathbb{R}^{r\times d_{\mathrm{in}}}$，$d_{\mathrm{in}},d_{\mathrm{out}}$分别为层输入与输出维数，$r$为内维数。以下固定相同缩放约定，并把缩放吸收入因子；完整训练机制及开销见第~\ref{sec:ket-transfer-added-parameters}~节。对任意可逆$\bm{R}\in\mathbb{R}^{r\times r}$，$(\bm{B}_k\bm{R})(\bm{R}^{-1}\bm{A}_k)$仍是同一增量，却有不同因子。不同任务还可能使用不同更新子空间，因此形状相同不足以使因子坐标可比。

KnOTS把对应问题放回完整增量：先形成各层$\Delta\bm{W}_k$，横向拼接，再作联合奇异值分解。这样用同一组输出侧基表示各任务增量，而不是分别相信各来源原有因子。对某一层可写为
\begin{equation}
  \begin{aligned}
    \relax[\Delta\bm{W}_1\ \cdots\ \Delta\bm{W}_K]
      &=\bm{U}\bm{\Sigma}[\bm{V}_1^\top\ \cdots\ \bm{V}_K^\top],\\
    \Delta\bm{W}_k&=\bm{U}\bm{\Sigma}\bm{V}_k^\top.
  \end{aligned}
  \label{eq:ket-integration-knots}
\end{equation}
其中$q$为拼接矩阵的秩，$\bm{U}\in\mathbb{R}^{d_{\mathrm{out}}\times q}$的列是共同正交基，$\bm{\Sigma}\in\mathbb{R}^{q\times q}$保存奇异值，$\bm{V}_k^\top\in\mathbb{R}^{q\times d_{\mathrm{in}}}$是第$k$个任务对应的块。不截断时等式精确成立；融合规则作用于这些共同坐标块，再乘回$\bm{U}\bm{\Sigma}$。这与融合后为压缩结果再分解不同。纯线性加和在精确分解下并未改变结果，坐标变化主要影响裁剪与符号选择等非线性操作。共同基提供运算的一致坐标，不会仅凭分解增加任务能力\scite{ket-stoica2025knots}

共同底座使增量引用同一版本、同一层和输入输出坐标，因而提供了直接对应的起点；只保存适配器时也必须保留准确底座版本。独立重训的底座不具备这一前提，不能把其适配器加到任意底座上。可能需要读取完整模型寻找层与激活对应，或者改用教师行为训练承接模型。深度、宽度或算子不同还会使逐参数运算没有定义，补零和同名层匹配不能证明语义兼容。对应可建立也仍需检查任务冲突，共同起点不是成功的充分条件。

\subsection{融合规则的构造}
\label{sec:ket-integration-merging-rule}

确定对应以后，要说明具体平均什么。将参与融合的参数块展平为$\bm{\theta}_k\in\mathbb{R}^{d_\theta}$，共同底座为$\bm{\theta}_0$，$d_\theta$为该参数块的坐标数。这里$\alpha_k$改指参数组合系数，不继承概率混合的非负与归一约束。直接加权完整参数，与加权相对底座的变化有何关系？在每项中加减底座，得到
\begin{equation}
  \sum_k\alpha_k\bm{\theta}_k
  =\left(\sum_k\alpha_k\right)\bm{\theta}_0+
    \sum_k\alpha_k(\bm{\theta}_k-\bm{\theta}_0).
  \label{eq:ket-integration-rule-equivalence}
\end{equation}
权重和为一时，直接平均与“底座加同系数增量”相同；否则完整参数加权还缩放了底座。这个恒等式只说明线性运算的关系，裁剪或因子相乘不能无条件照搬。下面分别讨论完整权重与增量的运算，系数如何由证据确定再由下一小节展开。

\subsubsection{完整参数的融合}
\label{sec:ket-integration-full-weights}

先固定同一花卉分类目标，并考虑同一个预训练模型经过不同微调得到的来源。Model Soups改变学习率、数据增强或训练随机性，保留这些微调权重，再对权重作平均。均匀方案对所有来源取$1/K$；贪心方案按验证准确率排序，从最佳来源开始依次试加，只在候选平均的验证准确率不下降时保留。也可学习插值系数。共同起点增加坐标直接对应的可能性，但贪心接纳仍是依据验证结果作出的选择，不保证独立测试也不下降\scite{ket-wortsman2022-soups}

等权平均不区分不同坐标对来源任务的敏感程度。Fisher merging用来源模型的对角Fisher信息构造局部二次近似，数值大的坐标对偏离来源值施加更强惩罚。具体地，$F_{k,i}\geq0$近似第$k$个来源在坐标$i$上的局部似然曲率；最小化$\sum_k F_{k,i}(\theta_i-\theta_{k,i})^2$，驻点给出
\[
  \theta_i^{\mathrm{F}}
  =\frac{\sum_k F_{k,i}\theta_{k,i}}
         {\sum_k F_{k,i}},
\]
分母为正时，这就是相应二次目标的最小点；分母为零则目标不约束该坐标，需预定回退值。Fisher的含义与估计区别已在前章的弹性权重巩固中解释，此处使用它来加权来源，而不是约束一次更新。统计必须由来源数据估计或随来源保存，对角近似也忽略相关方向。因此局部二次目标最优不等于真实多任务风险最优\scite{ket-matena2022-fisher}

RegMean换用线性层的输出偏差，要求融合层在各来源自己的输入上尽量复现原响应。为保持矩阵维度明确，此段采用行样本约定：$\bm{X}_k\in\mathbb{R}^{n_k\times d_{\mathrm{in}}}$，$\bm{W}_k\in\mathbb{R}^{d_{\mathrm{in}}\times d_{\mathrm{out}}}$，批输出为$\bm{X}_k\bm{W}_k$。取$n_k>0$，输入二阶矩阵为$\bm{G}_k=\bm{X}_k^\top\bm{X}_k/n_k$。最小化$\sum_k n_k^{-1}\lVert\bm{X}_k(\bm{W}-\bm{W}_k)\rVert_{\mathrm{F}}^2$，正规方程给出
\[
  \bm W^{\mathrm{R}}
  =\left(\sum_k\bm G_k\right)^{-1}
    \left(\sum_k\bm G_k\bm W_k\right),
\]
这里假设$\sum_k\bm{G}_k$可逆；病态或奇异时需要另定正则化与求解规则。相对于前文的列向量层权重，此处矩阵转置了，实际实现须保持约定一致。归一化使各来源贡献的是平均偏差，而不是让样本数大的来源天然占更大权重。原始数据可以由预先保存的二阶统计替代，但统计仍需来源输入产生；神经网络逐层应用这种规则也不保证整网输出偏差最小\scite{ket-jin2023-regmean}

这两种数据感知规则都在预测之前改变参数，不能视为来源预测的加权平均。参数平均运行$f(\bm{x};\sum_k\alpha_k\bm{\theta}_k)$，预测平均则先执行各个$f(\bm{x};\bm{\theta}_k)$。当函数对所融合参数满足相应线性关系时，平均次序才可交换；多层乘法与非线性激活通常破坏这种关系。下面固定来源与输入，仅改变平均的位置。

\begin{example}[非线性使两种平均不相等]
  \label{ex:ket-integration-nonlinear-average}
  取教学分数模型$f(x;w)=\operatorname{ReLU}(wx)$，来源参数$w_1=2,w_2=-1$，输入$x=1$，两权重都为$1/2$。两个来源分数为2和0，平均为1；参数平均为$\bar w=0.5$，再计算分数却得到$f(1;\bar w)=0.5$。在激活前平均与在激活后平均并不相等。

  将分数向量$(f(x;w),0)$归一化为二分类概率，第一坐标表示正类。两个来源正类概率为$e^2/(1+e^2)\approx0.8808$和$0.5$，平均约为$0.6904$；融合参数后的正类概率约为$0.6225$。先平均分数再做softmax得到约$0.7311$，又不同于平均概率。三种结果的来源、输入和权重完全相同，区别仅在运算顺序。
\end{example}

将物种模型与病害模型融合，还比同一分类目标下的多个解增加了任务差异。输出头的类别含义不同，就不能按同一坐标平均；可以只融合可对应的共享层，保留各任务自己的头。选择输出头可能仍需要任务身份，形成共享主干并没有解决未知任务的选择。Model Soups在同一目标上的结果不能直接证明任意技能可以融合。

任务变化以后，应在各来源任务上分别测风险，再按声明偏好汇总。若只运行一个共享主干和当前任务头，计算可能少于两个完整来源；保留多分支、重估归一化统计或校准则会增加其他成本。批归一化的运行均值与方差还依赖输入分布，不能未经核查就作为普通可训练权重一起平均。部署结构、状态处理和能力保留需要同一版本的记录。

Self-Soupervision进一步改变来源的训练目标，让共同底座在不同自监督目标或数据上继续训练。其方案既有之后再用标签微调的来源，也有直接融合自监督来源、用少量标签选系数的安排，以及无监督选择变体。它扩展了可融合来源的条件，但“继续训练不用标签”不等于“选择不用标签”，更不等于“评价不用标签”。若训练或选择访问了测试分布的无标签输入，就要声明直推式（亦称转导式）预测，即利用这批待预测输入共同准备模型，或声明相应适应协议；不能作为完全未见测试输入的普通归纳预测。这里仍按预印本的限定证据使用这项工作\scite{ket-fuller2026selfsoup}

\subsubsection{参数增量的融合}
\label{sec:ket-integration-deltas}

若目标是保留底座，再整合不同任务的微调变化，可以直接对变化量运算。\term{任务向量}（task vector）记录同一坐标下微调参数与共同底座的差；它是参数位移，不是一项已分离的纯语义技能。任务算术先作差，给各增量缩放系数，再加回底座：
\begin{equation}
  \bm{\tau}_k=\bm{\theta}_k-\bm{\theta}_0,\qquad
  \bm{\theta}_{\mathrm{merge}}(\bm{\lambda})
    =\bm{\theta}_0+\sum_{k=1}^{K}\lambda_k\bm{\tau}_k.
  \label{eq:ket-integration-task-arithmetic}
\end{equation}
其中$\bm{\tau}_k$为第$k$个任务向量，$\bm{\lambda}$汇集缩放系数。原始多任务形式通常让各$\lambda_k$相同，也可扩展为每来源独立系数。参数不是概率，系数无需非负或归一；负系数和较大尺度只是增加了可移动方向，不能由此推出能力可以相减或任意放大\scite{ket-ilharco2022taskarithmetic}

简单求和会让同一坐标的正负增量抵消，也会保留许多小幅变化。TIES-Merging改为先按绝对幅度保留给定比例的坐标，将其他值置零，得到$\widetilde{\bm{\tau}}_k$；再为每个坐标选择方向，只合并与该方向一致的非零增量。这里的方向由带幅度的和决定，合并采用同号平均：
\begin{equation}
  \begin{aligned}
    s_j&=\operatorname{sign}\!\left(\sum_k\widetilde{\tau}_{k,j}\right),\\
    I_j&=\{k:\widetilde{\tau}_{k,j}\neq0,\
                    \operatorname{sign}(\widetilde{\tau}_{k,j})=s_j\},\\
    \tau_{\mathrm{TIES},j}
      &=\frac{1}{|I_j|}\sum_{k\in I_j}\widetilde{\tau}_{k,j},\\
    \bm{\theta}_{\mathrm{merge}}
      &=\bm{\theta}_0+\lambda\bm{\tau}_{\mathrm{TIES}}.
  \end{aligned}
  \label{eq:ket-integration-ties}
\end{equation}
函数$\operatorname{sign}$返回正、负或零，$s_j$是坐标$j$的选举方向，$I_j$收集与它同号的来源，$\lambda$为最终缩放。本章约定$s_j=0$或$I_j$为空时输出零，幅度并列时按固定索引裁剪。方向不是按来源个数投票，平均分母也只计实际参与的来源；保留比例、并列规则与缩放都影响结果，应随版本记录\scite{ket-yadav2023ties}

用三维教学参数复算这一操作，底座与两来源为
\[
  \begin{aligned}
    \bm{\theta}_0&=(1,1,1)^\top,\\
    \bm{\theta}_1&=(1.4,0.8,1.01)^\top,\\
    \bm{\theta}_2&=(1.2,1.3,0.98)^\top.
  \end{aligned}
\]
按任务向量定义相减，得到
\[
  \begin{aligned}
    \bm{\tau}_1&=(0.4,-0.2,0.01)^\top,\\
    \bm{\tau}_2&=(0.2,0.3,-0.02)^\top.
  \end{aligned}
\]
每个来源保留两个最大幅度坐标后，增量为$(0.4,-0.2,0)^\top$和$(0.2,0.3,0)^\top$。前两维的和为$0.6,0.1$，都选正方向：第一维平均得$0.3$，第二维排除$-0.2$后只剩$0.3$。合并增量为$(0.3,0.3,0)^\top$，简单任务算术的和则为$(0.6,0.1,-0.01)^\top$。两者缩放都取一，加回底座分别得到$(1.3,1.3,1)^\top$与$(1.6,1.1,0.99)^\top$。这些计算确定参数不同，还未提供哪个任务风险更小的证据。

与按幅度选择不同，DARE以概率$p$随机丢弃增量坐标，再缩放保留值，之后可接任务算术或TIES。要让随机丢弃后仍保持原坐标的期望，就需用保留概率$1-p$作补偿。固定原增量，令$M_{k,j}\sim\operatorname{Bernoulli}(1-p)$，掩码为一时保留，且$0\leq p<1$，得到
\begin{equation}
  \begin{aligned}
    \widehat{\tau}_{k,j}&=\frac{M_{k,j}}{1-p}\tau_{k,j},\\
    \mathbb{E}_M[\widehat{\tau}_{k,j}]&=\tau_{k,j},&
    \operatorname{Var}_M(\widehat{\tau}_{k,j})
      &=\frac{p}{1-p}\tau_{k,j}^{2}.
  \end{aligned}
  \label{eq:ket-integration-dare}
\end{equation}
期望等式来自$\mathbb{E}_M[M_{k,j}]=1-p$，方差则由伯努利方差$p(1-p)$除以$(1-p)^2$得到。固定输入的单层线性更新也保持期望输出；非线性网络与任务损失却不能直接交换期望和计算。丢弃率接近一时，单坐标方差反而增大。实证能否接受这种随机变化取决于微调增量的尺度、冗余和丢弃率，不适用于任意独立训练模型\scite{ket-yu2024dare}

例如取$p=0.5$，前述$\bm{\tau}_1$的一次保留掩码为$(1,0,1)^\top$，结果为$(0.8,0,0.02)^\top$，并非原增量。无偏只描述所有掩码上的期望，不保证这一次保持行为。TIES按幅度和符号筛选，DARE随机筛选再补偿期望；二者可接在一起，却改变不同运算。

这些规则处理的都是坐标数值。隐藏单元重排能够保持函数，却改变逐坐标的正负对应，因此符号冲突不能直接当作语义冲突。同理，$\bm{\theta}_0-\lambda\bm{\tau}_k$只沿反向参数位移生成模型，不是精确删除某项知识。消除指定训练记录的影响需要第~\ref{sec:ket-update-unlearning-validation}~节的参照及\BookChapterRef{chap:trust-privacy}{卷3《基础、理论与可信性》第16章“隐私保护”}所述保证，不能由向量取负推出。

若希望整合结果仍保持低秩模块形态，可以在因子上运算，而不先对完整增量求和。LoRAHub对同一底座、同一插入层、相同秩和缩放约定的来源，分别加权$\bm{B}_k$和$\bm{A}_k$，再相乘形成模块：
\begin{equation}
  \begin{aligned}
    \bm{B}(\bm{w})&=\sum_k w_k\bm{B}_k,&
    \bm{A}(\bm{w})&=\sum_k w_k\bm{A}_k,\\
    \Delta\bm{W}(\bm{w})&=\bm{B}(\bm{w})\bm{A}(\bm{w}).
  \end{aligned}
  \label{eq:ket-integration-lorahub}
\end{equation}
其中$\bm{w}=(w_1,\ldots,w_K)^\top$为来源系数，可以为负且无需归一。这个结果是否等于完整增量加权，需要展开乘积，而不能仅看“加权”名称。两个来源时有
\begin{equation}
  \begin{aligned}
    \Delta\bm{W}(\bm{w})
      ={}&w_1^2\bm{B}_1\bm{A}_1+w_2^2\bm{B}_2\bm{A}_2\\
       &+w_1w_2(\bm{B}_1\bm{A}_2+\bm{B}_2\bm{A}_1).
  \end{aligned}
  \label{eq:ket-integration-lorahub-cross-terms}
\end{equation}
各来源自身的增量乘上系数平方，还出现$\bm{B}_1\bm{A}_2$和$\bm{B}_2\bm{A}_1$的交叉项，一般不同于$\sum_k w_k\Delta\bm{W}_k$。这些交叉项依赖来源因子的具体坐标，前文的可逆变换即使保持单个增量，也可能改变因子组合。原论文用$\bm{A}$命名输出侧因子，本章交换字母以统一维度，乘法次序保持不变\scite{ket-huang2024lorahub}

取秩一教学因子$\bm{B}_1=(1,0)^\top$、$\bm{A}_1=(1,0)$、$\bm{B}_2=(0,1)^\top$、$\bm{A}_2=(0,1)$，权重都为$1/2$。完整增量平均是对角矩阵$\operatorname{diag}(1/2,1/2)$，秩为二；因子平均后相乘，却得到所有元素均为$1/4$的矩阵，秩为一。非对角值来自交叉项，也说明保持因子内维数限制了可表达的秩。选定系数后只执行形成的模块，或将增量写回底座；AdapterFusion则保留多个适配器并组合各自激活，两者的运行路径不同。

\subsection{融合系数的确定}
\label{sec:ket-integration-coefficients}

同一融合规则会随系数产生多个候选，必须另有依据才能选择。预设等权系数不消耗目标选择标签，但也不能据此声称最优；有标签材料能直接比较任务损失，无标签材料只能提供熵或一致性等代理量。选择前应声明系数范围、访问权限和候选计算预算。测试标签一旦用于反复选系数，就已参加开发，不能再作为独立泛化证据。

LoRAHub在新任务的少量输入输出示例上评分，固定底座和来源因子，只搜索系数$\bm{w}$。每次先形成候选模块，再给目标输出计算负对数似然，并加入系数幅度惩罚；无梯度优化器依据候选评分决定下一组系数。原论文以Shiwa选择搜索策略，多数设置采用协方差矩阵自适应进化策略，即根据一批候选的评分更新采样分布。搜索达到预算后固定系数，预测时不再逐例重搜。“无梯度”省去的是搜索中的反向传播，不是示例、前向评分和搜索成本。

为了明确评分对象，令非空$S_u$为任务$u$的适应示例集，$\Omega$为预定系数范围，$p_{\bm{w}}$为底座加候选模块后的输出概率。取$\Omega$非空、闭且有界，并假设有限样本上的概率为正且随系数连续，目标就能取到最小值。用负对数概率使正确输出概率越大、损失越小，再以$\eta\geq0$限制系数幅度，理想搜索目标为
\begin{equation}
  \bm{w}_u^\star\in\argmin_{\bm{w}\in\Omega}
   \left\{-\frac{1}{|S_u|}
   \sum_{(\bm{x},y)\in S_u}\log p_{\bm{w}}(y\mid\bm{x})
   +\eta\lVert\bm{w}\rVert_1\right\}.
  \label{eq:ket-integration-lorahub-search}
\end{equation}
序列输出的对数概率按词元条件概率相加，具体实现的长度归一化与正则比例也应记录。本式概括负对数似然加幅度惩罚，不声称逐项复现实现细节。有限预算的无梯度搜索只能给出近似候选，不保证达到$\argmin$；增加来源会增大搜索维数，少量示例还会造成选择过拟合。最终任务评价应与$S_u$隔离。

若没有目标标签，AdaMerging用允许的无标签输入减小预测熵，寻找分布更集中的系数。任务级形式给每个来源任务向量一个$\lambda_k$；层级形式则给来源$k$在层$\ell$上的向量独立系数$\lambda_{k,\ell}$。固定底座和任务向量，只改变这些系数。用$\Omega_{\mathrm{A}}$记录实验允许的系数范围，一种带显式范围的目标写为
\begin{equation}
  \begin{aligned}
    \bm{\theta}_{\ell}(\bm{\lambda})
      &=\bm{\theta}_{0,\ell}+\sum_k\lambda_{k,\ell}\bm{\tau}_{k,\ell},\\
    \bm{\lambda}^{\star}
      &\in\argmin_{\bm{\lambda}\in\Omega_{\mathrm{A}}}\frac{1}{|U|}
        \sum_{\bm{x}\in U}H(p_{\bm{\lambda}}(\cdot\mid\bm{x})).
  \end{aligned}
  \label{eq:ket-integration-adamerging}
\end{equation}
$U$为非空无标签适应集，$H(p)=-\sum_y p(y)\log p(y)$为预测熵，$p_{\bm{\lambda}}$由候选参数和该任务接口给出。原研究从预设系数出发，用无标签测试样本的前向计算与熵梯度更新系数，按预算或停止规则结束。式中的$\Omega_{\mathrm{A}}$是为说明实验范围补入的记号，不把某一固定范围强加给所有实现。非空紧范围与连续概率足以保证最小点存在，但迭代仍不保证取得全局最小值；不具备存在条件时，只报告有限预算中得到的候选。任务级指每个来源向量共用一个系数，不是逐输入生成权重\scite{ket-yang2024adamerging}

熵衡量分布集中程度，不直接评价对错。恒把全部照片预测为月季，也能得到接近零的熵；非月季照片却都会出错。因而熵下降不能证明真实风险下降，还需真实标签评价，并检查是否塌缩为单类。用当前测试批次作$U$时属于直推式预测或测试时适应；要声称未来独立输入有效，就应冻结系数后另测。可见任务身份、适应轮数、持续更新和重置条件都应预先记录。

系数搜索还可能遇到无法同时满足的任务要求。设有$m$个目标任务，$R_j(\bm{\theta})$为任务$j$在其分布和损失下的平均风险，越小越好；风险向量为$(R_1(\bm{\theta}),\ldots,R_m(\bm{\theta}))$。一个可行参数若不存在另一可行参数在所有风险上不大、且至少一项更小，就是\term{帕累托最优}（Pareto optimality）。这表示不能不牺牲其他目标而继续改进某项风险，不表示每个任务都达到单独最优，更不表示算法已经找到了这种点。

Pareto Merging用偏好无关的共享参数和偏好相关的低秩修正构造候选族，让不同偏好得到不同折中。无数据时采用参数距离代理目标，有无标签输入时采用分任务熵目标；代理目标的帕累托性质不能直接等同于真实风险的全局最优。花卉系统可以分别绘制物种错误率与病害漏检率，再按预先声明的用途选择，而不把两者混成掩盖漏检退化的平均分\scite{ket-chen2025pareto}

这里的候选被已有来源与融合规则限制；第~\ref{sec:ket-sharing-directions}~节则在共同训练时协调各任务的参数下降方向。两者都处理多目标，但能改变的参数、可用训练证据和计算费用不同，不能借用相同术语交换保证。

表~\ref{tab:ket-integration-merging-operations}汇总运算对象与信息条件。完整参数、任务增量与低秩因子是不同运算对象，数据感知加权可作用于其中一些对象，对应调整又可先于多种规则。表的行因此允许交叉，不是一组互斥方法。第~\ref{sec:ket-integration-alignment}~节决定哪些坐标相容，不能由表中一种算术直接跳过对应检查。

来源存储、输入统计、校准材料以及反复构造候选和前向评分都属于整合成本。多轮吸收来源时，裁剪、匹配或重新选系数可能使顺序影响结果；应保存准确底座、来源、规则与系数，逐任务回归检查，并按有效维护要求选择回退版本。未来可塑性需要同预算的新任务学习实验，未见组合能力需要成分与关系分开的测试，都不是单份参数自动具有的性质。

Tam等把来源任务与未见组合任务分开，在图像分类、图像生成和跨语言任务的共同协议中比较融合方法。其预印本报告，来源任务成绩与未见任务成绩的关系随应用而变，增加来源时还可能出现来源任务保持下降、未见任务泛化上升的相反趋势。这项结果支持分开评价两者，不支持跨架构的固定排名\scite{ket-tam2024-realistic-merging}实际对照应包括底座、相关单来源、可执行的预测集成，以及数据和预算允许时的联合训练；来源能力保持也不应代替组合泛化测试。

\begin{table*}[htbp]
  \centering
  \small
  \setlength{\tabcolsep}{4pt}
  \renewcommand{\arraystretch}{1.2}
  \begin{tabularx}{\linewidth}{@{}>{\raggedright\arraybackslash}p{24mm}
      >{\raggedright\arraybackslash}X
      >{\raggedright\arraybackslash}X
      >{\raggedright\arraybackslash}X@{}}
    \toprule
    操作 & 运算对象与可执行条件 & 系数所用信息 & 部署与主要风险 \\
    \midrule
    完整参数平均 & 全部权重或共享层；要求层、坐标和输出含义对应 & 预设系数，或在独立开发材料上选择、学习 & 通常单模型，也可多头；平均模型可能落入高损失区域 \\
    \addlinespace
    数据感知加权 & 参数坐标或线性层；Fisher merging用局部曲率，RegMean用输入二阶矩 & 来源统计或生成统计的输入；仍可另选超参数 & 单模型；统计失配、对角近似和病态矩阵限制结论 \\
    \addlinespace
    任务增量融合 & 相对同一准确底座的位移；可接裁剪、符号选择或随机稀疏 & 固定缩放、标签选择或无标签适应 & 底座加增量；任务冲突、过大位移或随机波动可能损害行为 \\
    \addlinespace
    低秩因子组合 & 同底座、插入层、秩与缩放约定；分别加权因子再相乘 & LoRAHub用少量任务示例作无梯度搜索 & 一个模块或写回底座；交叉项、因子坐标与样本过拟合影响结果 \\
    \midrule
    对应调整（前置步骤） & 置换、特征匹配或共同子空间；结构异构时仍可能无法运算 & 权重匹配只读参数，激活匹配另需输入 & 可后接融合或保留分支；对应成立不保证任务能力兼容 \\
    \bottomrule
  \end{tabularx}
  \caption{参数融合的运算对象、所需信息和部署结果。上部各行允许交叉，下部的对应调整可以先于多种规则执行。无标签只限制标签访问，不表示没有输入、任务身份或适应计算。比较来源保留、未见组合和资源节省时，均须以最终实际模型而非操作名称为准。}
  \label{tab:ket-integration-merging-operations}
\end{table*}

\section{基于知识蒸馏的知识整合}
\label{sec:ket-integration-distillation}

两个花卉分类器的结构不同，参数可能无法对应，输出却仍可以使用同一类别顺序。这时可以查询来源，将其回答作为训练目标，让新模型承接这些行为。提供监督的来源称为教师，接受监督的模型称为学生。教师也可由组合系统或融合模型担任；学生可以更小，也可以改用另一种结构。可行性由监督接口、输入覆盖和学生能够表达的行为决定，不由模型大小关系决定。

\subsection{教师与蒸馏输入的准备}
\label{sec:ket-integration-teachers-inputs}

首先确认教师实际提供什么。分类接口只返回类别时，学生只能获得硬答案；返回概率或分数时才有相应软信息；能读取内部激活时才可设置表示约束。若两教师都回答物种问题，应按第~\ref{sec:ket-integration-prediction}~节统一类别、尺度和聚合规则。若分别回答物种与病害，就要保留任务条件或独立监督头，不能把不同问题混成一个向量。

教师之间也可能共同出错或给出冲突答案。蒸馏前应固定模仿对象：等权集成、条件集成、任务专用教师，或融合后的模型。先集成再蒸馏时，要同时测单教师、教师集成和学生，区分来源聚合的收益与学生训练的收益。第~\ref{chap:ket-reuse-selection}~章的可靠性和收益判断仍适用，多个教师不意味着增加了同等数量的独立信息。

教师只能在输入被查询时提供约束，因此还要确定查询覆盖。原训练输入接近教师熟悉的区域，却可能包含记忆性响应；另采无标签输入能补充条件，也可能超出教师可靠范围；合成输入能补充难采样场景，但生成偏差和重复模板会限制覆盖。若只有月季和百合查询，没有兰花输入，蒸馏损失即使为零，也没有约束学生在兰花上的回答。查询区拟合不能证明全部来源能力保持。

训练时的信息还可能多于部署时。广义蒸馏让教师使用学生不可见的特权信息，例如教师同时读取照片和检测报告，学生只有照片；教师为配对的照片生成监督\scite{ket-lopezpaz2016generalized}配对必须对应同一植株与观测时刻，前章式~\eqref{eq:ket-update-crossmodal}已说明表示对应的约束。若照片无法区分病原而报告可以，学生仍不能对每例恢复只在报告中出现的差异。

为明确这种不可见信息的限度，固定学生照片$\bm{x}$，将教师额外观测记为$\bm{x}^{+}$，教师给类别$y$的软概率为$q_y(\bm{x},\bm{x}^{+})$。学生只能给同一$\bm{x}$分配一个概率向量。若允许选择任意概率分布，最小化额外观测条件下的期望交叉熵，最优预测是
\begin{equation}
  p^\star(y\mid\bm{x})
   =\mathbb{E}\!\left[q_y(\bm{x},\bm{x}^{+})\mid\bm{x}\right].
  \label{eq:ket-integration-privileged-limit}
\end{equation}
期望目标可写为$-\sum_y\mathbb{E}[q_y\mid\bm{x}]\log p(y\mid\bm{x})$，与条件平均教师分布的熵相减，剩下非负KL散度，故两分布相等时取到最小值。零概率采用通常的极限约定。这个理想解是可见信息下的平均教师判断，不是对隐藏报告的逐例恢复；有限学生容量和有限样本还可能使它无法达到。

\subsection{蒸馏约束的构造}
\label{sec:ket-integration-distillation-losses}

查询接口确定以后，需要选择匹配哪一部分行为。输出约束比较学生与教师的最终答案或分布，表示约束比较内部计算；两者可以共同使用。也可从这些量派生样本间关系矩阵、位置注意力图或导数约束，但必须另外定义比较对象与访问条件。本节集中建立输出和表示这两种基本约束，不把派生信号视为互斥的第三类载体。

\subsubsection{基于输出行为的蒸馏}
\label{sec:ket-integration-output-distillation}

教师判成月季时，仍可能认为它更像百合而不像菊花。硬类别丢掉这些次要倾向，概率目标则保留了它们。经典蒸馏用softmax温度调整概率的集中程度，使非最大类别获得更明显的监督贡献\scite{hinton2015distilling}

固定同一照片及$c$个类别顺序，教师和学生分数为$\bm{a}^{\mathrm{T}},\bm{a}^{\mathrm{S}}\in\mathbb{R}^{c}$，上标$\mathrm{T},\mathrm{S}$标记角色。将分数除以共同温度$T>0$再归一化，得到
\begin{equation}
  q_j^{(T)}=\frac{\exp(a_j^{\mathrm{T}}/T)}
                    {\sum_{v=1}^{c}\exp(a_v^{\mathrm{T}}/T)},\qquad
  p_j^{(T)}=\frac{\exp(a_j^{\mathrm{S}}/T)}
                    {\sum_{v=1}^{c}\exp(a_v^{\mathrm{S}}/T)}.
  \label{eq:ket-integration-temperature}
\end{equation}
温度上标$(T)$不是学习迭代。温度一恢复普通softmax；分数不全相同时，升温使分布趋向均匀。要让学生接近冻结的教师分布，可以最小化教师概率加权的负对数学生概率，也就是交叉熵：
\begin{equation}
  L_{\mathrm{out}}
  =-\sum_j q_j^{(T)}\log p_j^{(T)},\qquad
  D_{\mathrm{KL}}(\bm{q}^{(T)}\Vert\bm{p}^{(T)})
  =L_{\mathrm{out}}-H(\bm{q}^{(T)}).
  \label{eq:ket-integration-output-loss}
\end{equation}
其中$D_{\mathrm{KL}}(\bm{q}\Vert\bm{p})=\sum_jq_j\log(q_j/p_j)$，衡量教师分布到学生分布的偏离，方向在此固定为教师在前。它非负且相等时为零；交叉熵只比它多一个冻结教师的熵，因而对学生有相同梯度。多教师先平均分数再软化，通常不同于各自软化后平均概率。只取得概率时，须使用声明的概率目标，或依据前章已给出的完整概率与温度转换条件重温度化，不能假装取得了原始分数。

\Needspace{10\baselineskip}
\begin{example}[三个类别携带的信息]
  \label{ex:ket-integration-soft-labels}
  一张教学照片的教师分数按月季、百合、菊花排列为$(2,1,0)$。用上述归一化分别计算温度一与温度二，得到
  \[
    \begin{aligned}
      T=1:\quad &(0.6652,0.2447,0.0900),\\
      T=2:\quad &(0.5065,0.3072,0.1863).
    \end{aligned}
  \]
  硬目标$(1,0,0)$仅保留月季，而软目标还要求百合概率高于菊花。另一个教师若给$(2,0,1)$，最大类别仍是月季，次要类别却反序，因此软监督不同。升温增加了次要类别的概率贡献，不说明这些倾向一定符合真值。
\end{example}

改变温度还会改变训练步的尺度。为判断软目标项怎样与真实标签项共同使用，对学生分数求导：
\[
  \frac{\partial\log p_j^{(T)}}{\partial a_v^{\mathrm{S}}}
    =\frac{\mathbf{1}\{j=v\}-p_v^{(T)}}{T}.
\]
将该导数代入交叉熵，并利用教师概率之和为一，得到
\begin{equation}
  \frac{\partial L_{\mathrm{out}}}{\partial a_v^{\mathrm{S}}}
   =-\frac{1}{T}\sum_jq_j^{(T)}
      \bigl(\mathbf{1}\{j=v\}-p_v^{(T)}\bigr)
   =\frac{p_v^{(T)}-q_v^{(T)}}{T}.
  \label{eq:ket-integration-temperature-gradient}
\end{equation}
导数显式含$1/T$，但概率差也会随温度缩小。固定有限分数，记$\bar a^{\mathrm{S}}$、$\bar a^{\mathrm{T}}$为两者各自类别分数的均值；温度远大于中心化分数的幅度时，用指数一阶展开并归一化，可得
\begin{equation}
  \begin{aligned}
    p_v^{(T)}
      &=\frac{1}{c}+\frac{a_v^{\mathrm{S}}-\bar a^{\mathrm{S}}}{cT}
          +O(T^{-2}),\\
    \frac{\partial L_{\mathrm{out}}}{\partial a_v^{\mathrm{S}}}
      &=\frac{(a_v^{\mathrm{S}}-\bar a^{\mathrm{S}})
              -(a_v^{\mathrm{T}}-\bar a^{\mathrm{T}})}
              {cT^2}+O(T^{-3}).
  \end{aligned}
  \label{eq:ket-integration-temperature-square}
\end{equation}
学生与教师的概率都偏离均匀分布约$1/T$，再乘导数中的$1/T$，使高温梯度约按$1/T^2$缩小。常用$T^2L_{\mathrm{out}}$或$T^2D_{\mathrm{KL}}$就是补偿这一尺度，与温度一的真实标签项保持大致可比；它不是任意温度下的精确梯度等价。例~\ref{ex:ket-integration-soft-labels}~中学生初始分数全零、$T=2$时，月季梯度约为$(1/3-0.5065)/2=-0.0866$，乘温度平方后约为$-0.3463$。梯度下降都会提高月季分数，但步幅不同。

输出还可以从类别扩展为动作。Policy Distillation比较硬动作、动作价值回归与动作分布匹配；动作价值描述在当前状态选择该动作后预期的折扣回报，可经softmax形成教师目标。动作编号、可用动作和状态观测必须对应，不能把同编号的不同动作视为同一信号。该研究还使用低温突出价值较高的动作，所以分类中软化分布的选择不能机械套到策略蒸馏\scite{ket-rusu2015policy}

动作选择会改变随后看到的状态，教师轨迹不必覆盖学生执行时到达的全部状态；这项新增条件不能由分类输入覆盖直接满足，完整交互分析可参考\BookVolumeRef{vol:reinforcement-learning}{卷7《学习范式：强化学习》}。分类也要保留任务条件：同一照片分别问物种与病害时需要不同答案，若不给学生任务信息，监督冲突仍在。蒸馏损失下降仅证明接近所选信号，真实任务准确率仍需标签检验。

\subsubsection{基于中间表示的蒸馏}
\label{sec:ket-integration-feature-distillation}

只限制最终类别，还没有告诉学生中间层应保留哪些可供后续计算的信息。FitNets选取教师隐藏层输出作为提示，让学生某个中间层学习预测这个提示。教师层较宽、学生层较窄时维度不同，先用学生侧可学习映射把表示转成相容形式，再计算距离\scite{ket-romero2015fitnets}

教师第$\ell_{\mathrm{T}}$层表示为$\bm{z}^{\mathrm{T}}\in\mathbb{R}^{d_{\mathrm{T}}}$，学生第$\ell_{\mathrm{S}}$层为$\bm{z}^{\mathrm{S}}\in\mathbb{R}^{d_{\mathrm{S}}}$，映射$r_{\bm{\phi}}:\mathbb{R}^{d_{\mathrm{S}}}\to\mathbb{R}^{d_{\mathrm{T}}}$。在固定输入上，平方残差惩罚各对应坐标的偏离，提示损失为
\begin{equation}
  L_{\mathrm{rep}}(\bm{x})
   =\frac{1}{2}\left\|
      r_{\bm{\phi}}(\bm{z}^{\mathrm{S}}(\bm{x}))
        -\bm{z}^{\mathrm{T}}(\bm{x})
     \right\|_2^2.
  \label{eq:ket-integration-feature-loss}
\end{equation}
这里$\bm{\phi}$和生成学生表示的参数参与学习，教师表示固定。系数$1/2$简化梯度，不改变平方损失的最小点。卷积提示还需对应空间位置、通道和感受野；映射可以调整形状，但插值或聚合丢失的信息须说明。距离的数值也依赖尺度，不能只凭向量等长就比较。

用三维教师提示与二维学生表示计算一次教学更新：$\bm{z}^{\mathrm{T}}=(1,0,2)^\top$、$\bm{z}^{\mathrm{S}}=(0.5,1)^\top$，取线性映射
\[
  r_{\bm{\phi}}(\bm{z}^{\mathrm{S}})=\bm{R}\bm{z}^{\mathrm{S}},
  \qquad
  \bm{R}=\begin{bmatrix}2&0\\0&1\\0&2\end{bmatrix}.
\]
同一输入分别经过教师和学生，再经映射得到$(1,1,2)^\top$，残差为$(0,1,0)^\top$，故损失为$1/2$。对$\bm{R}$求导得到残差与学生表示的外积，对学生表示求导得到$\bm{R}^{\top}$乘残差；提示梯度因此只进入映射和学生此前各层。暂时固定学生、以学习率$0.1$更新映射，第二行由$(0,1)$变成$(-0.05,0.9)$，映射后的第二坐标为$0.875$，损失为$0.3828125$。这展示了怎样学习维度对应，不是把短向量补零。

FitNets原始安排先训练学生至提示层及映射，再用输出蒸馏训练完整学生。将两者同时优化是另一种可选安排，不能混作原算法。提示映射只为训练提供比较接口时，部署可以移除它，最终预测仍沿学生自己的输出路径；若其他方案把映射接入预测路径，则不能同样移除。

直接逐坐标匹配要求内部坐标对应，学习映射则允许学生采用能预测教师提示的另一种表示。第~\ref{sec:ket-integration-alignment}~节已用置换说明内部坐标不是语义本身，过强的多层匹配可能限制学生自己的表示。另一方面，容量很大的映射可能记住有限训练集的对应，低提示损失仍不证明新输入上可读出相同信息。应检验独立数据上的提示与最终任务效果，并按实际部署路径测学生；跨结构、跨模态还须重新核对位置和可观测信息。

\subsection{承接模型的训练}
\label{sec:ket-integration-student-training}

输出和表示约束确定后，要把它们接到可执行训练中。令$S_{\mathrm{L}}$为真实标签训练集，$U$为蒸馏输入集，二者可以重叠，但不含最终测试材料。学生参数为$\bm{\theta}$，提示映射参数为$\bm{\phi}$；教师和聚合规则固定。真实标签纠正任务答案，软目标限制输出偏离，提示限制中间残差，各项先按样本平均再赋非负权重，一种联合目标为
\begin{equation}
  \begin{aligned}
    L(\bm{\theta},\bm{\phi})={}&
      \frac{\lambda_{\mathrm{L}}}{|S_{\mathrm{L}}|}
      \sum_{(\bm{x},y)\in S_{\mathrm{L}}}
        -\log p_{\bm{\theta}}(y\mid\bm{x};1)\\
      &+\frac{\lambda_{\mathrm{D}}T^2}{|U|}
      \sum_{\bm{x}\in U}
        D_{\mathrm{KL}}\!\left(
          \bm{q}^{(T)}(\bm{x})\Vert
          \bm{p}_{\bm{\theta}}^{(T)}(\bm{x})\right)\\
      &+\frac{\lambda_{\mathrm{R}}}{|U|}
      \sum_{\bm{x}\in U}L_{\mathrm{rep}}(\bm{x};\bm{\theta},\bm{\phi}).
  \end{aligned}
  \label{eq:ket-integration-joint-distillation}
\end{equation}
$\lambda_{\mathrm{L}},\lambda_{\mathrm{D}},\lambda_{\mathrm{R}}$分别控制三项权重，$p_{\bm{\theta}}(y\mid\bm{x};1)$为温度一的学生概率。软目标项在教师与学生上用同一$T$，并采用此前解释的$T^2$尺度补偿。启用的项要求相应集合非空；没有标签就去掉第一项，没有内部提示就去掉第三项，而不是对空集作除法。温度、权重和容量在允许的验证材料上选定。联合目标只定义折中，学生容量不足或信号矛盾时，不保证所有项都能达到各自最小值。

\begin{algorithm}[从冻结教师到可部署学生]
  \label{alg:ket-integration-distillation-loop}
  \begin{algorithmic}[1]
    \Require 固定版本的教师与相容接口，可用的$S_{\mathrm{L}}$、$U$，独立验证集，损失设置、优化器与训练预算
    \Ensure 所选学生参数、部署配置和验证记录
    \State 初始化学生与所需提示映射，冻结教师及其依赖路径
    \For{学习步骤$t=0,1,\ldots$，直到预算耗尽或满足停止条件}
      \State 按预定任务与类别比例，从各启用目标的训练集合抽取小批次
      \State 查询冻结教师或读取兼容缓存，取得软目标与可用提示
      \State 沿学生路径计算所需概率、提示映射与各项损失
      \State 按式~\eqref{eq:ket-integration-joint-distillation}汇总启用项
      \State 用预定优化器只更新学生及提示映射
      \State 按验证任务损失和保留要求检查停止条件，保存候选检查点
    \EndFor
    \State 按预定验证规则选择并冻结学生，移除仅训练需要的教师与映射
    \State 用未参与选择的测试材料检查任务、来源保持、组合能力与实际部署成本
  \end{algorithmic}
\end{algorithm}

算法~\ref{alg:ket-integration-distillation-loop}中的采样决定各任务实际贡献，常见任务比例过大可能削弱少见任务的保留。教师缓存必须绑定版本、预处理与温度；增强改变照片后，原缓存未必仍对应当前输入。停止可依据预算、验证损失停止改善或保留要求恶化，条件应事先确定。部署一般采用温度一，额外概率校准则单独记录其参数和选择材料。

比较效果时，固定任务、输入范围和测试样本，报告教师与蒸馏学生；有真实标签时，再加入同结构、同允许预算而只用标签训练的学生。没有训练标签时，这一对照不可执行，应在相同可用输入和预算下比较其他可行基线。较小容量增加逼近困难，异构结构增加提示匹配与求解条件；远程教师还受查询次数、输出精度和内部不可见限制。匹配损失并没有改变这些容量、求解与访问条件。

底座升级将参数不对应的问题具体化。Trans-LoRA把旧底座加旧适配器作为教师，在新底座上训练新适配器，而不是把旧因子直接装上去。论文的迁移阶段保留五个任务种子示例，并使用原任务阶段训练的判别器：新底座生成候选输入，判别器筛选，旧模型为保留输入提供输出，新适配器学习这些输出。原完整数据迁移时不再读取，但种子、生成、筛选、查询与训练仍有数据和计算成本。论文的近乎无原数据设定不保证任意任务无损迁移\scite{ket-wang2024translora}

长期使用还要检查压缩误差是否累积。Progress \& Compress保留知识库网络和活动分支：当前任务先训练活动分支，并读取冻结知识库的表示；压缩时反过来以活动分支为教师，将新行为蒸馏回知识库，同时用在线弹性权重巩固限制重要参数变化。两部分交替承担当前学习与历史压缩，结构保持固定规模，而非不断增加任务列\scite{ket-schwarz2018progress}

这种教师若读取旧知识库，压缩阶段必须固定它所依赖的旧版本或缓存，否则更新知识库也会悄然移动监督目标。参数保留不能补回未查询的区域，也不能无限扩大容量。若每轮只查询常见花卉，少见物种可能逐次偏离；学生随后再作教师，还会复制已有错误。应逐任务复测保留、通用能力与同预算后续学习，并累计教师存储、生成、查询和重复训练费用。持续验证由第~\ref{sec:ket-update-retention-validation}~节补充，压缩算子与部署实现仍归模型卷和\BookPartRef{part:systems}{卷12《系统》第1部分}。

\section{本章小结}
\label{sec:ket-integration-summary}

回到已有两个分类器和若干属性模块的花卉系统，多来源并不直接意味着多种能力都能使用。组合保留各来源，却要统一答案含义或中间接口，并承担实际调用成本；固定与条件只说明规则是否依赖任务或输入，不说明它是否经过训练。互补错误可使平均纠正部分答案，共同错误和错误路由也会保留下来。

若希望只运行一个模型，参数融合需要先建立对应，再选择运算与系数。相同函数也可能因隐藏单元重排而不能直接平均，低秩因子相乘又会引入完整增量平均没有的交叉项。蒸馏无需相同参数坐标，却需要可靠教师、可比监督、查询覆盖和可训练学生。二次近似最优、熵下降或蒸馏拟合都只回答其对应目标，不能替代真实任务能力的证据。

因此，选择应同时服从信息与资源条件：能够持续执行来源时可组合，参数可对应时可试融合，更换结构或压缩行为时可训练学生，也可将前两者再蒸馏。最终分别检验来源任务保持、未见组合和后续学习，记录准备、部署与长期维护成本。如何为新目标准备并形成行为由第~\ref{chap:ket-transfer-adaptation}~章继续展开，多任务共同形成知识由第~\ref{chap:ket-sharing}~章展开，各项主张的统一评价协议见第~\ref{chap:ket-evaluation}~章。
